\ifdefined\gkver\else\def\gkver{student}\fi
\documentclass[\gkver]{gaokaozhenti}
\begin{document}

\chapter{2009年浙江理综}
%% paperType: 理综物理部分

\begin{enumerate}
\item
%% number: 14
%% paperName: 2009年浙江理综第 14 题 6 分
%% typeId: 1
%% type: 单选题
%% chapter: 力物体的平衡
%% point: 三力平衡
%% method: 无
%% score: 6
%% degree: 650
%% duplicateId: 0
%% body:
如图所示，质量为 $m$ 的等边三棱柱静止在水平放置的斜面上。已知三棱柱与斜面之的动摩擦因数为 $\mu$，斜面的倾角为 $30^\circ$，则斜面对三棱柱的支持力与摩擦力的大小分别为\xzanswer{A}

\onepicture[w=5.2cm]{figs/14.png}

\fourchoices[answer=A]
{$\frac{\sqrt{3}}{2}mg$ 和 $\frac{1}{2}mg$}
{$mg$ 和 $\frac{\sqrt{3}}{2}mg$}
{$\frac{1}{2}mg$ 和 $\frac{1}{2}\mu mg$}
{$\frac{\sqrt{3}}{2}mg$ 和 $\frac{\sqrt{3}}{2}\mu mg$}

%% answer: A
%% memo:
\memoanswer{
受力如图所示，$F_{N}=mg\cos30^\circ=\frac{\sqrt{3}}{2}mg$，$F_{f}=mg\sin30^\circ=\frac{1}{2}mg$。

\onepicture[w=3.6cm]{figs/14_an.png}
}

\item
%% number: 15
%% paperName: 2009年浙江理综第 15 题 6 分
%% typeId: 1
%% type: 单选题
%% chapter: 原子和原子核
%% point: 原子核式结构
%% method: 无
%% score: 6
%% degree: 850
%% duplicateId: 0
%% body:
氦原子核由两个质子与两个中子组成，这两个质子之间存在着万有引力、库仑力和核力，则 3 种力从大到小的排列顺序是\xzanswer{D}

\fourchoices[answer=D]
{核力、万有引力、库仑力}
{万有引力、库仑力、核力}
{库仑力、万有引力、核力}
{核力、库仑力、万有引力}

%% answer: D
%% memo:
\memoanswer{
核力是强力，它能将核子束缚在原子核内。万有引力最弱，研究核子间相互作用时万有引力可以忽略。
}

\item
%% number: 16
%% paperName: 2009年浙江理综第 16 题 6 分
%% typeId: 1
%% type: 单选题
%% chapter: 电场
%% point: 库仑定律
%% method: 无
%% score: 6
%% degree: 650
%% duplicateId: 0
%% body:
如图所示，在光滑绝缘水平面上放置 3 个电荷量均为 $q$（$q>0$）的相同小球，小球之间用劲度系数均为 $k_{0}$ 的轻质弹簧绝缘连接。当 3 个小球处在静止状态时，每根弹簧长度为 $l_{0}$。已知静电力常量为 $k$，若不考虑弹簧的静电感应，则每根弹簧的原长为\xzanswer{C}

\onepicture[w=6.2cm]{figs/16.png}

\fourchoices[answer=C]
{$l+\frac{5kq^{2}}{2k_{0}l^{2}}$}
{$l-\frac{kq^{2}}{k_{0}l^{2}}$}
{$l-\frac{5kq^{2}}{4k_{0}l^{2}}$}
{$l-\frac{5kq^{2}}{2k_{0}l^{2}}$}

%% answer: C
%% memo:
\memoanswer{
第三个小球受三个力的作用，它们的关系是 $k_{0}x=k\frac{q^{2}}{l^{2}}+k\frac{q^{2}}{(2l)^{2}}$，得 $x=\frac{5kq^{2}}{4k_{0}l^{2}}$；$l_{0}=l-x=l-\frac{5kq^{2}}{4k_{0}l^{2}}$。
}

\item
%% number: 17
%% paperName: 2009年浙江理综第 17 题 6 分
%% typeId: 1
%% type: 单选题
%% chapter: 电磁感应
%% point: 楞次定律
%% method: 无
%% score: 6
%% degree: 450
%% duplicateId: 0
%% body:
如图所示，在磁感应强度大小为 $B$、方向竖直向上的匀强磁场中，有一质量为 $m$、阻值为 $R$ 的闭合矩形金属线框 abcd 用绝缘轻质细杆悬挂在 $O$ 点，并可绕 $O$ 点摆动。金属线框从右侧某一位置静止开始释放，在摆动到左侧最高点的过程中，细杆和金属框平面始终处于同一平面，且垂直纸面。则线框中应电流的方向是\xzanswer{B}

\onepicture[w=4.6cm]{figs/17.png}

\fourchoices[answer=B]
{$a\to b\to c\to d\to a$}
{$d\to c\to b\to a\to d$}
{先是 $d\to c\to b\to a\to d$，后是 $a\to b\to c\to d\to a$}
{先是 $a\to b\to c\to d\to a$，后是 $d\to c\to b\to a\to d$}

%% answer: B
%% memo:
\memoanswer{
由楞次定律，一开始磁通量减小，后来磁通量增大，由“增反”“减同”可知电流方向是 $d\to c\to b\to a\to d$。
}

\item
%% number: 18
%% paperName: 2009年浙江理综第 18 题 6 分
%% typeId: 2
%% type: 多选题
%% chapter: 光学
%% point: 全反射
%% method: 无
%% score: 6
%% degree: 650
%% duplicateId: 0
%% body:
如图所示，有一束平行于等边三棱镜截面 ABC 的单色光从空气射向 E 点，并偏折到 F 点。已知入射方向与边 AB 的夹角为 $\theta=30^\circ$，E、F 分别为边 BC 的中点，则\xzanswer{AC}

\onepicture[w=5.0cm]{figs/18.png}

\fourchoices[answer=AC]
{该棱镜的折射率为 $\sqrt{3}$}
{光在 F 点发生全反射}
{光从空气进入棱镜，波长变小}
{从 F 点出射的光束与入射到 E 点的光束平行}

%% answer: AC
%% memo:
\memoanswer{
在 E 点作出法结线可知入射角为 $60^\circ$，折射角为 $30^\circ$，折射率为 $\sqrt{3}$；由光路的可逆性可知，在 BC 边上的入射角小于临界角，不会发生全反射，B 错；由公式 $\lambda_{\text{介}}=\frac{\lambda_{\text{空气}}}{n}$，可知 C 对；三棱镜两次折射使得光线都向底边偏折，不会与入射到 E 点的光束平行，故 D 错。
}

\item
%% number: 19
%% paperName: 2009年浙江理综第 19 题 6 分
%% typeId: 2
%% type: 多选题
%% chapter: 曲线运动
%% point: 万有引力定律
%% method: 无
%% score: 6
%% degree: 650
%% duplicateId: 0
%% body:
在讨论地球潮汐成因时，地球绕太阳运行轨道与月球绕地球运行轨道可视为圆轨道。已知太阳质量约为月球质量的 $2.7\times10^{7}$ 倍，地球绕太阳运行的轨道半径约为月球绕地球运行的轨道半径的 400 倍。关于太阳和月球对地球上相同质量海水的引力，以下说法正确的是\xzanswer{AD}

\fourchoices[answer=AD]
{太阳引力远大于月球引力}
{太阳引力与月球引力相差不大}
{月球对不同区域海水的吸引力大小相等}
{月球对不同区域海水的吸引力大小有差异}

%% answer: AD
%% memo:
\memoanswer{
$\frac{F_{\text{太阳}}}{F_{\text{月}}}=\frac{M_{\text{太阳}}}{M_{\text{月}}}\cdot\frac{R_{\text{月}}^{2}}{R_{\text{太阳}}^{2}}$，代入数据可知，太阳的引力远大于月球的引力；由于月心到不同区域海水的距离不同，所以引力大小有差异。
}

\item
%% number: 20
%% paperName: 2009年浙江理综第 20 题 6 分
%% typeId: 2
%% type: 多选题
%% chapter: 电场
%% point: 带电粒子的直线运动
%% method: 无
%% score: 6
%% degree: 450
%% duplicateId: 0
%% body:
空间存在匀强电场，有一电荷量 $q$（$q>0$），质量 $m$ 的粒子从 O 点以速率 $v_{0}$ 射入电场，运动到 A 点时速率为 $2v_{0}$。现有另一电荷为 $q$、质量 $m$ 的粒子以速率 $2v_{0}$ 仍从 O 点射入该电场，运动到 B 点时速率为 $3v_{0}$。若忽略重力的影响，则\xzanswer{AD}

\fourchoices[answer=AD]
{在 O、A、B 三点中，B 点电势最高}
{在 O、A、B 三点中，A 点电势最高}
{OA 间的电势差比 BO 间的电势差大}
{OA 间的电势差比 BA 间的电势差小}

%% answer: AD
%% memo:
\memoanswer{
正电荷由 O 到 A，动能变大，电场力做正功，电势能减小，电势也减小，O 点电势较高；负电荷从 O 到 B 速度增大，电场力也做正功，电势能减小，电势升高，B 点电势比 O 点高。所以 B 点最高，A 对；$U_{OA}=\frac{W_{OA}}{q}=\frac{\frac{1}{2}m(2v_{0})^{2}-\frac{1}{2}m(v_{0})^{2}}{q}=\frac{3mv_{0}^{2}}{2q}$，$U_{OB}=\frac{W_{OB}}{-q}=\frac{\frac{1}{2}m(3v_{0})^{2}-\frac{1}{2}m(2v_{0})^{2}}{-q}=\frac{5mv_{0}^{2}}{-2q}$，故 D 对。
}

\item
%% number: 21
%% paperName: 2009年浙江理综第 21 题 6 分
%% typeId: 1
%% type: 单选题
%% chapter: 机械波
%% point: 波与振动图象
%% method: 无
%% score: 6
%% degree: 450
%% duplicateId: 0
%% body:
一列波长大于 $1\Um$ 的横波沿着 $x$ 轴正方向传播，处在 $x_{1}=1\Um$ 和 $x_{2}=2\Um$ 的两质点 A、B 的振动图像如图所示。由此可知\xzanswer{A}

\onepicture[w=6.2cm]{figs/21.png}

\fourchoices[answer=A]
{波长为 $\frac{4}{3}\Um$}
{波速为 $1\Ums$}
{3 s 末 A、B 两质点的位移相同}
{1 s 末 A 点的振动速度大于 B 点的振动速度}

%% answer: A
%% memo:
\memoanswer{
$\Delta x=x_{2}-x_{1}=1\Um$，由于波沿 $x$ 正方向传播，所以 A 先振动，又由于波长大于 $1\Um$，所以 $\Delta t=3\Us=\frac{3}{4}T$，所以 $\Delta x=\frac{3}{4}\lambda$，$\lambda=\frac{4}{3}\Um$，A 对，波速 $v=\frac{\Delta x}{\Delta t}=\frac{1}{3}\Ums$，B 错；由振动图像知，在 3 s 末，A、B 两质点的位移不相同，C 错；1 s 末 A 点速度为零，B 点速度最大，D 错。
}

\item
%% number: 22
%% paperName: 2009年浙江理综第 22 题 9 分
%% typeId: 4
%% type: 实验题
%% chapter: 机械振动
%% point: 单摆测重力加速度
%% method: 无
%% score: 9
%% degree: 450
%% duplicateId: 0
%% body:
\begin{enumerate}
	\item
	如图甲所示，在“描绘小灯泡的伏安特性曲线”实验中，同组同学已经完成部分导线的连接，请你在实物接线图中完成余下导线的连接。
	\onepicture[num=甲, w=6.0cm]{figs/22a.png}
	\item
	某同学从标称为“$220\UV$ $25\UW$”“$220\UV$ $300\UW$”“$220\UV$ $500\UW$”的 3 只灯泡中任选一只，正确使用多用电表测量灯泡阻值如图乙所示。该灯泡的阻值是 \tkanswer[1.6cm]{} $\UO$，标称的额定功率为 \tkanswer[1.6cm]{} $\UW$。
	\onepicture[num=乙, w=7.2cm]{figs/22b.png}
	\item
	在“探究单摆周期与摆长的关系”实验中，两位同学用游标卡尺测量小球的直径如图甲、乙所示。测量方法正确的是 \tkanswer[1.6cm]{}（选填“甲”或“乙”）。
	\twopicture[numstyle=chinese, w=6.5cm]{figs/22c.png}{figs/22d.png}
	\item
	实验时，若摆球在垂直纸面的平面内摆动，为了将人工记录振动次数改为自动记录振动次数，在摆球运动最低点的左、右两侧分别放置一激光光源与光敏电阻，如图甲所示。光敏电阻与某一自动记录仪相连，该仪器显示的光敏电阻阻值 $R$ 随时间 $t$ 变化图线如图乙所示，则该单摆的振动周期为 \tkanswer[1.6cm]{}。若保持悬点到小球顶点的绳长不变，改用直径是原小球直径 2 倍的另一小球进行实验，则该单摆的周期将 \tkanswer[1.6cm]{}（填“变大”、“不变”或“变小”），图乙中的 $\Delta t$ 将 \tkanswer[1.6cm]{}（填“变大”、“不变”或“变小”）。
	\onepicture[num=甲, w=4.2cm]{figs/22e.png}
	\onepicture[num=乙, w=9.2cm]{figs/22f.png}
\end{enumerate}

%% answer:
\jdanswer{
\begin{enumerate}
	\item
	见解析图。
	\item
	$160\UO$；$25\UW$。
	\item
	乙。
	\item
	$2t_{0}$；变大；变大。
\end{enumerate}
}

%% memo:
\memoanswer{
Ⅰ．本题考查多用电表和电路设计

（1）如图所示。实验要求电压从零开始调整，滑动变阻器采用分压式接法。

\onepicture[w=8.6cm]{figs/22_an.png}

（2）$16\times10=160\UO$，由于灯泡正常工作时的电阻为 $R=\frac{U^{2}}{P}$ 比不工作时的电阻大的多，而“$220\UV$ $25\UW$”“$220\UV$ $300\UW$”“$220\UV$ $500\UW$”正常工作时的电阻分别为 $1936\UO$、$161\UO$ 和 $98\UO$，故为 $25\UW$。

Ⅱ．本题考查游标卡尺和单摆的知识

（1）应将待测物体正确地放在测脚中如乙图；（2）单摆 1 个周期遮光两次；单摆周期与小球质量、大小无关，但若改用直径变为原小球直径的 2 倍，周期变大，但遮光时间 $\Delta t$ 变大。
}

\item
%% number: 23
%% paperName: 2009年浙江理综第 23 题 14 分
%% typeId: 6
%% type: 计算题
%% chapter: 电场
%% point: 带电粒子的直线运动
%% method: 无
%% score: 14
%% degree: 450
%% duplicateId: 0
%% body:
如图所示，相距为 $d$ 的平行金属板 A、B 竖直放置，在两板之间水平放置一绝缘平板。有一质量 $m$、电荷量 $q$（$q>0$）的小物块在与金属板 A 相距 $l$ 处静止。若某一时刻在金属板 A、B 间加一电压 $U_{AB}=-\frac{3\mu mgd}{2q}$，小物块与金属板只发生了一次碰撞，碰撞后电荷量变为 $-\frac{1}{2}q$，并以与碰前大小相等的速度反方向弹回。已知小物块与绝缘平板间的动摩擦因数为 $\mu$，若不计小物块几何量对电场的影响和碰撞时间。则：

\begin{enumerate}
	\item
	小物块与金属板 A 碰撞前瞬间的速度大小是多少？
	\item
	小物块碰撞后经过多长时间停止运动？停在何位置？
\end{enumerate}
\onepicture[w=5.0cm, align=right]{figs/23.png}

%% answer:
\jdanswer{
\begin{enumerate}
	\item
	$\sqrt{\mu gl}$
	\item
	$4\sqrt{\frac{l}{\mu g}}$，停在 $2l$ 处或距离 B 板为 $d-2l$
\end{enumerate}
}

%% memo:
\memoanswer{
（1）加电压后，B 极板电势高于 A 板，小物块在电场力作用与摩擦力共同作用下向 A 板做匀加速直线运动。电场强度为 $E=\frac{U_{BA}}{d}$，小物块所受的电场力与摩擦力方向相反，则合外力为 $F_{\text{合}}=qE-\mu mg$，故小物块运动的加速度为 $a_{1}=\frac{F_{\text{合}}}{m}=\frac{qU-\mu mgd}{md}=\frac{1}{2}\mu g$。设小物块与 A 板相碰时的速度为 $v_{1}$，由 $v_{1}^{2}=2a_{1}l$，解得 $v_{1}=\sqrt{\mu gl}$。

（2）小物块与 A 板相碰后以 $v_{1}$ 大小相等的速度反弹，因为电荷量及电性改变，电场力大小与方向发生变化，摩擦力的方向发生改变，小物块所受的合外力大小为 $F_{\text{合}}=\mu mg-\frac{qE}{2}$，加速度大小为 $a_{2}=\frac{F_{\text{合}}}{m}=\frac{1}{4}\mu g$。设小物块碰后到停止的时间为 $t$，注意到末速度为零，有 $0-v_{1}=-a_{2}t$，解得 $t=\frac{v_{1}}{a_{2}}=4\sqrt{\frac{l}{\mu g}}$。设小物块碰后停止时距离为 $x$，注意到末速度为零，有 $0-v_{1}^{2}=-2a_{2}x$，则 $x=\frac{v_{1}^{2}}{2a_{2}}=2l$，或距离 B 板为 $d-2l$。
}

\item
%% number: 24
%% paperName: 2009年浙江理综第 24 题 18 分
%% typeId: 6
%% type: 计算题
%% chapter: 功和能
%% point: 交通工具启动
%% method: 无
%% score: 18
%% degree: 450
%% duplicateId: 0
%% body:
某校物理兴趣小组决定举行遥控塞车比赛。比赛路径如图所示，赛车从起点 A 出发，沿水平直线轨道运动 $L$ 后，出 B 点进入半径为 $R$ 的光滑竖直圆轨道，离开竖直圆轨道后继续在光滑平直轨道上运动到 C 点，并能越过壕沟。已知赛车质量 $m=0.1\Ukg$，通电后以额定功率 $P=1.5\UW$ 工作，进入竖直圆轨道前受到的阻值为 $0.3\UN$，随后在运动中受到的阻力均可不计。图中 $L=10.00\Um$，$R=0.32\Um$，$h=1.25\Um$，$S=1.50\Um$。问：要使赛车完成比赛，电动机至少工作多长时间？（取 $g=10\Umsq$）

\onepicture[w=8.0cm, align=right]{figs/24.png}

%% answer:
\jdanswer{$2.53\Us$}

%% memo:
\memoanswer{
本题考查平抛、圆周运动和功能关系。

设赛车越过壕沟需要的最小速度为 $v_{1}$，由平抛运动的规律 $S=v_{1}t$，$h=\frac{1}{2}gt^{2}$，解得 $v_{1}=S\sqrt{\frac{R}{2h}}=3\Ums$。

设赛车恰好越过圆轨道，对应圆轨道最高点的速度为 $v_{2}$，最低点的速度为 $v_{3}$，由牛顿第二定律及机械能守恒定律 $mg=m\frac{v_{2}^{2}}{R}$，$\frac{1}{2}mv_{3}^{2}=\frac{1}{2}mv_{2}^{2}+mg(2R)$，解得 $v_{3}=\sqrt{5gh}=4\Ums$。

通过分析比较，赛车要完成比赛，在进入圆轨道前的速度最小应该是 $v_{\min}=4\Ums$。

设电动机工作时间至少为 $t$，根据功能原理 $Pt-fL=\frac{1}{2}mv_{\min}^{2}$，由此可得 $t=2.53\Us$。
}

\item
%% number: 25
%% paperName: 2009年浙江理综第 25 题 22 分
%% typeId: 6
%% type: 计算题
%% chapter: 磁场
%% point: 带电粒子在磁场中的运动
%% method: 无
%% score: 22
%% degree: 250
%% duplicateId: 0
%% body:
如图所示，$x$ 轴正方向水平向右，$y$ 轴正方向竖直向上。在 $xOy$ 平面内与 $y$ 轴平行的匀强电场，在半径为 $R$ 的圆内还有与 $xOy$ 平面垂直的匀强磁场。在圆的左边放置一带电微粒发射装置，它沿 $x$ 轴正方向发射出一束具有相同质量 $m$、电荷量 $q$（$q>0$）和初速度 $v$ 的带电微粒。发射时，这束带电微粒分布在 $0<y<2R$ 的区间内。已知重力加速度大小为 $g$。

\begin{enumerate}
	\item
	从 A 点射出的带电微粒平行于 $x$ 轴从 C 点进入有磁场区域，并从坐标原点 O 沿 $y$ 轴负方向离开，求电场强度和磁感应强度的大小与方向。
	\item
	请指出这束带电微粒与 $x$ 轴相交的区域，并说明理由。
	\item
	在这束带电磁微粒初速度变为 $2v$，那么它们与 $x$ 轴相交的区域又在哪里？并说明理由。
\end{enumerate}
\onepicture[w=7.0cm, align=right]{figs/25.png}

%% answer:
\jdanswer{
\begin{enumerate}
	\item
	$E=\frac{mg}{q}$，方向沿 $y$ 轴正方向；$B=\frac{mv}{qR}$，方向垂直于纸面向外
	\item
	这束带电微粒都通过坐标原点。
	\item
	这束带电微粒与 $x$ 轴相交的区域是 $x>0$。
\end{enumerate}
}

%% memo:
\memoanswer{
本题考查带电粒子在复合场中的运动。

带电粒子平行于 $x$ 轴从 C 点进入磁场，说明带电微粒所受重力和电场力平衡。设电场强度大小为 $E$，由 $mg=qE$，可得 $E=\frac{mg}{q}$，方向沿 $y$ 轴正方向。

带电微粒进入磁场后，将做圆周运动。且 $r=R$，如图（a）所示，设磁感应强度大小为 $B$。由 $qvB=\frac{mv^{2}}{R}$，得 $B=\frac{mv}{qR}$，方向垂直于纸面向外。

\onepicture[w=5.4cm]{figs/25_an1.png}

（2）这束带电微粒都通过坐标原点。

方法一：从任一点 P 水平进入磁场的带电微粒在磁场中做半径为 $R$ 的匀速圆周运动，其圆心位于其正下方的 Q 点，如图 b 所示，这束带电微粒进入磁场后的圆心轨迹是如图 b 的虚线半圆，此圆的圆心是坐标原点。

方法二：从任一点 P 水平进入磁场的带电微粒在磁场中做半径为 $R$ 的匀速圆周运动。如图 b 所示，设 P 点与 $O'$ 点的连线与 $y$ 轴的夹角为 $\theta$，其圆心 Q 的坐标为 $(-R\sin\theta, R\cos\theta)$，圆周运动轨迹方程为 $(x+R\sin\theta)^{2}+(y-R\cos\theta)^{2}=R^{2}$，得 $x=0$，$y=0$，或 $x=-R\sin\theta$，$y=R(1+\cos\theta)$。

\onepicture[w=5.4cm]{figs/25_an2.png}

（3）这束带电微粒与 $x$ 轴相交的区域是 $x>0$。

带电微粒在磁场中经过一段半径为 $r'$ 的圆弧运动后，将在 $y$ 轴的右方（$x>0$）的区域离开磁场并做匀速直线运动，如图 c 所示。靠近 M 点发射出来的带电微粒在突出磁场后会射向 $x$ 轴正方向的无穷远处，靠近 N 点发射出来的带电微粒会在靠近原点之处穿出磁场。所以，这束带电微粒与 $x$ 轴相交的区域范围是 $x>0$。

\onepicture[w=7.4cm]{figs/25_an3.png}
}

\end{enumerate}

\end{document}
